% Imporditud: tools/import_eksperimendid.py
% Allikas: Eesti füüsikaolümpiaadi eksperimendi ülesannete
%   kogu (Rael Kalda, 2023), ülesanne Ü161, lahendus L161.
\setAuthor{EFO žürii}
\setRound{lõppvoor}
\setYear{2000}
\setNumber{G E2}
\setDifficulty{5}
\setTopic{elektriahelad}
\setVahendid{nelja väljundklemmiga kinnine karp, oommeeter (multimeeter). Takistuse mõõtmisel on soovitav kasutada oommeetri mõõtepiirkonda ``2000 $\Omega$''}
\setPunktid{15}
\setAllikas{Eksperimendi kogu Ü161, lahendus lk 121}
\setLahendusseis{toores}

\prob{Elektriskeem karbis}
Kinnises karbis on neljast detailist koosnev elektriskeem. Joonistage see skeem ja leidke detailide takistused.

\hint

\solu
Kinnises kastis on järgmine skeem (vt. joon.). Numbrid klemmide juures vastavad karbil olevatele numbritele. Arvulisi väärtusi: $R_1$ = (100 $\pm$ 10) $\Omega$ ja $R_2$ = (50 $\pm$ 5) $\Omega$. Dioodide takistused on vahemikus 1000 $\Omega$ kuni 2000 $\Omega$.

Mõõtmine: Kõigepealt kontrollitakse multimeetriga vooluallikate olemasolu. Et antud juhul need puuduvad, siis mõõdetakse takistust iga klemmi vahel, vahetades multimeetri polaarsust. Andmete põhjal koostatakse võrrandisüsteem, mille lahendamise tulemusel saadakse elementide takistused.

Hindamine: Takistuse mõõtmine iga võimaliku klemmipaari vahel mõlema polaarsuse jaoks (4 p). Võrrandite koostamine ja lahendamine (8 p). Õige elektriskeem (3 p).
\probend
